\section{Background and Related Work}
\label{sec:related-work}

\subsection{Concept Erasure and Its Reference}

Post-training concept erasure aims to suppress a model's knowledge of a target
concept while preserving behavior unrelated to that concept. More broadly,
machine unlearning is commonly framed as approximating the model that would
have resulted had the data to be forgotten not contributed to training
\citep{wichert2024rethinking}. This framing suggests two distinct evaluation
questions. First, does an intervention reduce target knowledge while
preserving other capabilities? Second, does the resulting model resemble one
trained without the target learning signal?

Common evaluations compare an edited model with its original version, measuring
target suppression and capability preservation but not similarity to a model
trained without the target data \citep{yao2024machineunlearning,li2024wmdp}.
A matched training intervention is therefore needed to evaluate the latter.

\subsection{Concept Exclusion with LMEnt}

LMEnt provides a Wikipedia-based pretraining corpus with entity annotations,
an index for retrieving chunks associated with particular entities, and
language models trained on the corpus \citep{gottesman2025lment}. Its entity
index makes it possible to define a concept-level training intervention:
chunks linked to a selected set of entities can be prevented from contributing
to the training loss.

We refer to a model trained under such an intervention as a
\emph{concept-excluded twin}. A twin is trained under otherwise matched
conditions but receives no loss on the selected concept-associated chunks. It
therefore provides a reference for the behavioral effect of excluding a
defined portion of the training signal.

Concept exclusion does not guarantee that a model contains no target
knowledge. The entity-based exclusion set may miss indirect or unlinked
references while also covering text about neighboring topics. Consequently,
a twin may retain target knowledge and differ from the full model outside the
target concept. It therefore represents the effect of the specified
loss-exclusion procedure, not a perfectly concept-free model.

Prior work has used controlled changes to pretraining data to study how
knowledge and capabilities are acquired. Hubble inserts selected text and
compares models to examine memorization \citep{wei2025hubble}, whereas Deep
Ignorance removes dual-use material before pretraining and tests whether the
corresponding capabilities remain absent, including after adversarial
fine-tuning \citep{obrien2025deepignorance}. We use a related approach for a
different purpose: for each concept, we train a matched twin with the same
initialization, batch order, and training configuration as the full model, but
mask the loss on entity-linked chunks. This twin represents the outcome of our
specified concept-exclusion procedure and provides a reference for evaluating
post-training erasure.

\subsection{Post-Training Concept-Erasure Methods}

We study three post-training methods that intervene at different locations in
the model. EMBER edits input-token embeddings rather than internal transformer
layers \citep{suslik2026ember}. It applies sparse matrix factorization to
embeddings obtained from target and neutral sentences, identifies factors
associated with the target concept, and subtracts scaled contributions of
those factors from selected token embeddings. The remaining model parameters
are left unchanged.

Representation Misdirection for Unlearning (RMU) modifies internal model
representations \citep{li2024wmdp}. It fine-tunes selected transformer weights
so that representations of forget-set inputs approach a random control
direction, while a retain objective encourages representations on non-target
inputs to remain close to those of the original model. RMU therefore combines
a target-directed intervention with an explicit preservation objective.

Semi-nonnegative matrix factorization (SNMF) identifies interpretable features
from groups of co-activated MLP neurons \citep{shafran2026snmf}. In the
erasure adaptation evaluated by EMBER, concept-associated features are mapped
to directions that can be removed through edits to MLP weights
\citep{suslik2026ember}. Because EMBER acts on token embeddings whereas RMU
and the SNMF-based method act on internal weights, EMBER can also be combined
with either MLP-level method in a staged intervention.

\subsection{Evaluating Erasure Against Concept Exclusion}

Prior evaluations of concept erasure have primarily measured whether an edit
reduces target knowledge while preserving unrelated behavior. EMBER evaluates
target efficacy, related-domain specificity, general-capability preservation,
output coherence, and resistance to relearning \citep{suslik2026ember}. RMU
similarly measures suppression on WMDP alongside performance on MMLU and
MT-Bench \citep{li2024wmdp}. These evaluations establish whether an edit is
effective and selective, but they do not compare the edited model with a
matched model trained without the target learning signal. We address this gap
by constructing a concept-excluded twin for each target and measuring whether
post-training erasure moves the model toward that reference. In addition to
suppression and preservation, we compare correct-answer negative
log-likelihood and teacher-forced full-vocabulary distributions on the same
sentence completions.
